\chapter{Requirements \& Design}
\label{ch:requirements-design}

% Structure mirrors the host lab's thesis format (7 sections):
%   3.1 State Of The Art Informed Requirements
%   3.2 User Centered Requirements
%   3.3 Design Requirements Consolidation
%   3.4 User Experience (UX) Design
%   3.5 User Interface (UI) Design
%   3.6 Moderator Agent Design   (= "Assistant Design" in the host-lab format)
%   3.7 Use Case

This chapter presents the requirements that drove the design of AIutami
and the resulting design choices. The requirements are first derived
from the literature gaps surfaced in Chapter~\ref{ch:state-of-the-art}
(Section~\ref{sec:rd-sota-req}) and from user-centred enquiry with
target participants and the supervising lab
(Section~\ref{sec:rd-user-req}). Section~\ref{sec:rd-consolidation}
consolidates the two streams into a unified requirement list.
Sections~\ref{sec:rd-ux},~\ref{sec:rd-ui},
and~\ref{sec:rd-moderator-agent} describe respectively the user
experience, the user interface, and the moderator agent design.
Section~\ref{sec:rd-use-case} closes the chapter with a representative
use case that walks through a full group session.

\section{State Of The Art Informed Requirements}
\label{sec:rd-sota-req}
% TODO: requirements derived from the gaps surfaced in Chapter 2.
% Examples:
%  - Support N participants in real-time speech (multi-party, not dyadic)
%  - Explicit turn-taking (mitigate multi-party scalability problem)
%  - Pluggable task layer (NASA Moon, Lost at Sea, Murder Mystery, Generic)
%  - Comparable control arm (moderator OFF) for empirical evaluation

\section{User Centered Requirements}
\label{sec:rd-user-req}
% TODO: requirements gathered from informal feedback with the host lab and
% from the pilot participants (tutors). Examples:
%  - Mobile-friendly UI (participants often join from phones)
%  - Robust rejoin after disconnect (mobile networks are unstable)
%  - Visible turn indicator (no doubt about who is speaking)
%  - GDPR-compliant data collection (Italian university study)

\section{Design Requirements Consolidation}
\label{sec:rd-consolidation}
% TODO: merge SOTA-derived + user-derived into a single numbered list
% (R1, R2, ...). Cite each one back to its source section.

\section{User Experience (UX) Design}
\label{sec:rd-ux}
% TODO: describe the user journey:
% Lobby -> Individual ranking (NASA/LAS only) -> Group discussion (ACTIVE)
% -> Group ranking submission -> Conclusion -> Report.
% Mention design choices: avatar of moderator, ready-to-conclude button,
% participation indicators, etc.

\section{User Interface (UI) Design}
\label{sec:rd-ui}
% TODO: describe the screens (lobby, ranking, session, conclusion, report).
% Include wireframes / screenshots. Discuss accessibility choices.

\section{Moderator Agent Design}
\label{sec:rd-moderator-agent}
% TODO: design rationale for the AI moderator:
%  - Goals: neutrality, minimal intrusion, fairness
%  - Trigger taxonomy: NO_PUSH, INACTIVE_VOICE, TIMER_25/30, soft (LLM)
%  - Intervention categories: questioning, summarising, redirecting, time-keeping
%  - Forced conclusion recap
%  - Prompt engineering decisions

\section{Use Case}
\label{sec:rd-use-case}
% TODO: walk through one representative session (e.g. Lost at Sea, 4 participants,
% moderator ON). Step-by-step narrative covering lobby join, intro, individual
% ranking, group discussion with one or two AI interventions, ready-to-conclude
% flow, and report generation.
